\chapter{Những bài học từ người bình thường}
\markboth{Những bài học từ người bình thường}{Lessons from Ordinary People}


\section{Người lao động chăm chỉ}
\begin{truyen}{Người lao động chăm chỉ}{Life lesson}
\chuhoa{B}{ác công nhân dậy từ bốn giờ, về tối mịt. Bác nói: làm để con đi học. Người lao động chăm chỉ.}
\textit{(The worker got up at four, came back at dark. He said: I work so my kids can study. The hardworking laborer.)}
\end{truyen}

\begin{baihoc}
Chăm chỉ là phẩm giá.
\textit{(Hard work is dignity.)}
\end{baihoc}

\begin{ghinhoanh}
\textbf{Short English to remember:}

Hard work is dignity.

\end{ghinhoanh}

\begin{tuhocnhanh}
1. Câu chuyện này gợi cho bạn suy nghĩ gì? \textit{What does this story make you think?}

2. Bạn sẽ nhớ bài học nào nhất? \textit{Which lesson will you remember most?}
\end{tuhocnhanh}
\ngancach


\section{Người mẹ hy sinh}
\begin{truyen}{Người mẹ hy sinh}{Life lesson}
\chuhoa{M}{ẹ nhịn ăn để con no, nhịn mặc để con đẹp. Người mẹ hy sinh.}
\textit{(Mother went hungry so her children could eat, wore old clothes so they could look nice. The sacrificing mother.)}
\end{truyen}

\begin{baihoc}
Hy sinh thầm lặng là tình yêu.
\textit{(Silent sacrifice is love.)}
\end{baihoc}

\begin{ghinhoanh}
\textbf{Short English to remember:}

Silent sacrifice is love.

\end{ghinhoanh}

\begin{tuhocnhanh}
1. Câu chuyện này gợi cho bạn suy nghĩ gì? \textit{What does this story make you think?}

2. Bạn sẽ nhớ bài học nào nhất? \textit{Which lesson will you remember most?}
\end{tuhocnhanh}
\ngancach


\section{Người cha ít nói}
\begin{truyen}{Người cha ít nói}{Life lesson}
\chuhoa{C}{ha không nói nhiều, nhưng mỗi lần con cần tiền học hay cần lời khuyên, cha có. Người cha ít nói.}
\textit{(Father didn't speak much, but whenever his child needed tuition or advice, he was there. The quiet father.)}
\end{truyen}

\begin{baihoc}
Hành động nói thay lời.
\textit{(Actions speak instead of words.)}
\end{baihoc}

\begin{ghinhoanh}
\textbf{Short English to remember:}

Actions speak instead of words.

\end{ghinhoanh}

\begin{tuhocnhanh}
1. Câu chuyện này gợi cho bạn suy nghĩ gì? \textit{What does this story make you think?}

2. Bạn sẽ nhớ bài học nào nhất? \textit{Which lesson will you remember most?}
\end{tuhocnhanh}
\ngancach


\section{Người bán hàng kiên nhẫn}
\begin{truyen}{Người bán hàng kiên nhẫn}{Life lesson}
\chuhoa{C}{hị bán rau từ sáng sớm, khách mặc cả, chị vẫn cười. Người bán hàng kiên nhẫn.}
\textit{(She sold vegetables from dawn, customers haggled, she still smiled. The patient vendor.)}
\end{truyen}

\begin{baihoc}
Kiên nhẫn nuôi sống và nuôi tâm.
\textit{(Patience feeds the body and the soul.)}
\end{baihoc}

\begin{ghinhoanh}
\textbf{Short English to remember:}

Patience feeds the body and the soul.

\end{ghinhoanh}

\begin{tuhocnhanh}
1. Câu chuyện này gợi cho bạn suy nghĩ gì? \textit{What does this story make you think?}

2. Bạn sẽ nhớ bài học nào nhất? \textit{Which lesson will you remember most?}
\end{tuhocnhanh}
\ngancach


\section{Người thầy tận tụy}
\begin{truyen}{Người thầy tận tụy}{Life lesson}
\chuhoa{T}{hầy dạy thêm không lấy tiền cho học trò nghèo. Người thầy tận tụy.}
\textit{(The teacher tutored for free for poor students. The devoted teacher.)}
\end{truyen}

\begin{baihoc}
Tận tụy là món quà.
\textit{(Devotion is a gift.)}
\end{baihoc}

\begin{ghinhoanh}
\textbf{Short English to remember:}

Devotion is a gift.

\end{ghinhoanh}

\begin{tuhocnhanh}
1. Câu chuyện này gợi cho bạn suy nghĩ gì? \textit{What does this story make you think?}

2. Bạn sẽ nhớ bài học nào nhất? \textit{Which lesson will you remember most?}
\end{tuhocnhanh}
\ngancach


\section{Người bạn trung thành}
\begin{truyen}{Người bạn trung thành}{Life lesson}
\chuhoa{K}{hi anh bệnh nặng, bạn ấy đến mỗi ngày. Người bạn trung thành.}
\textit{(When he was seriously ill, that friend came every day. The loyal friend.)}
\end{truyen}

\begin{baihoc}
Trung thành là vàng.
\textit{(Loyalty is gold.)}
\end{baihoc}

\begin{ghinhoanh}
\textbf{Short English to remember:}

Loyalty is gold.

\end{ghinhoanh}

\begin{tuhocnhanh}
1. Câu chuyện này gợi cho bạn suy nghĩ gì? \textit{What does this story make you think?}

2. Bạn sẽ nhớ bài học nào nhất? \textit{Which lesson will you remember most?}
\end{tuhocnhanh}
\ngancach


\section{Người già từng trải}
\begin{truyen}{Người già từng trải}{Life lesson}
\chuhoa{Ô}{ng kể chuyện đời, anh nghe. Ông nói: cháu đừng lặp sai của ông. Người già từng trải.}
\textit{(The old man told his life story, he listened. The man said: don't repeat my mistakes. The experienced elder.)}
\end{truyen}

\begin{baihoc}
Kinh nghiệm của người già là kho báu.
\textit{(Elders' experience is treasure.)}
\end{baihoc}

\begin{ghinhoanh}
\textbf{Short English to remember:}

Elders' experience is treasure.

\end{ghinhoanh}

\begin{tuhocnhanh}
1. Câu chuyện này gợi cho bạn suy nghĩ gì? \textit{What does this story make you think?}

2. Bạn sẽ nhớ bài học nào nhất? \textit{Which lesson will you remember most?}
\end{tuhocnhanh}
\ngancach


\section{Người trẻ đầy hy vọng}
\begin{truyen}{Người trẻ đầy hy vọng}{Life lesson}
\chuhoa{C}{ô gái mới ra trường nói: em sẽ cố. Ánh mắt sáng. Người trẻ đầy hy vọng.}
\textit{(The new graduate said: I'll try. Her eyes were bright. The hopeful young.)}
\end{truyen}

\begin{baihoc}
Hy vọng là sức mạnh.
\textit{(Hope is strength.)}
\end{baihoc}

\begin{ghinhoanh}
\textbf{Short English to remember:}

Hope is strength.

\end{ghinhoanh}

\begin{tuhocnhanh}
1. Câu chuyện này gợi cho bạn suy nghĩ gì? \textit{What does this story make you think?}

2. Bạn sẽ nhớ bài học nào nhất? \textit{Which lesson will you remember most?}
\end{tuhocnhanh}
\ngancach


\section{Người thất bại nhưng không bỏ cuộc}
\begin{truyen}{Người thất bại nhưng không bỏ cuộc}{Life lesson}
\chuhoa{A}{nh phá sản hai lần, vẫn đứng dậy. Người thất bại nhưng không bỏ cuộc.}
\textit{(He went bankrupt twice, still got up. The one who failed but didn't quit.)}
\end{truyen}

\begin{baihoc}
Không bỏ cuộc mới là thắng.
\textit{(Not quitting is winning.)}
\end{baihoc}

\begin{ghinhoanh}
\textbf{Short English to remember:}

Not quitting is winning.

\end{ghinhoanh}

\begin{tuhocnhanh}
1. Câu chuyện này gợi cho bạn suy nghĩ gì? \textit{What does this story make you think?}

2. Bạn sẽ nhớ bài học nào nhất? \textit{Which lesson will you remember most?}
\end{tuhocnhanh}
\ngancach


\section{Người sống giản dị}
\begin{truyen}{Người sống giản dị}{Life lesson}
\chuhoa{B}{à không cần nhà lớn, xe đẹp. Bà nói: đủ ăn, con cháu sum vầy là đủ. Người sống giản dị.}
\textit{(She didn't need a big house or a nice car. She said: enough to eat, family together is enough. The simple liver.)}
\end{truyen}

\begin{baihoc}
Giản dị là biết đủ.
\textit{(Simplicity is knowing enough.)}
\end{baihoc}

\begin{ghinhoanh}
\textbf{Short English to remember:}

Simplicity is knowing enough.

\end{ghinhoanh}

\begin{tuhocnhanh}
1. Câu chuyện này gợi cho bạn suy nghĩ gì? \textit{What does this story make you think?}

2. Bạn sẽ nhớ bài học nào nhất? \textit{Which lesson will you remember most?}
\end{tuhocnhanh}
\ngancach
